% Annexe de méthode commune, écrite le 16/11/2026 ; base 2025 le 17/11/2026 ; coupée le 17/11/2026
% (synthese/CONSIGNES_COUPE.md) : reçoit le bilan détaillé du 1 % (ann:meth_haut, tab:ens_haut), sorti de la synthèse.
% Sources : revue/methode_perte_seche.md, revue/methode_v2.md, synthese/perte_seche_autres.py,
% synthese/synthese_ensemble.py. Les calibrages propres à chaque réforme sont dans les annexes de partie (ann:ircsg,
% ann:foncier, ann:is, ann:pfu) et ne sont pas repris ici.

\subsection{Une seule mesure de l'efficacité : la perte sèche généralisée}\label{ann:meth_ps}

La note mesure de combien chaque réforme réduit la perte sèche par rapport à aujourd'hui, à recette égale, avec la
même formule pour les quatre (Harberger, 1964 ; Hines, 1999) :
\[
\Delta PS \;=\; -\sum_j \big(\bar T_j - s_j\,\tau_{k(j)}\big)\,\Delta Q_j ,\qquad
\Delta Q_j = \varepsilon_j\,Q_j\,\Delta\ln(1-t_j).
\]
L'indice $j$ parcourt les marges sur lesquelles la réforme agit ; $Q_j$ est la quantité taxée (revenu d'activité,
ventes de logements, stock d'équipements, revenus du capital déclarés) ; $t_j$, le taux réformé auquel elle réagit ;
$\bar T_j$, le coin total sur la marge, moyenne de ses valeurs avant et après ; $s_j$, la part de la réaction qui se
déplace vers une autre assiette taxée au taux $\tau_k$ (Chetty, 2009). Sur une marge déjà taxée, un point de plus coûte
au premier ordre le coin existant multiplié par la quantité perdue : beaucoup là où le coin est lourd, comme au
voisinage du Smic.

\paragraph{Le coin total} compte tout ce qu'un euro d'activité perdu fait perdre à la collectivité : pour l'IR-CSG,
l'impôt fusionné, les cotisations qui n'ouvrent pas de droits et le retrait de la prime d'activité (\ref{ann:ircsg}) ;
pour l'immobilier, le droit de mutation sur les ventes et l'impôt sur le bâti rapporté au coût d'usage sur la
construction, le sol n'ayant pas de marge ; pour l'IS, le coin sur le coût d'usage des équipements (Hall-Jorgenson).

\paragraph{Le réel et le déclaratif.} Seule la baisse réelle de l'activité ou de l'investissement entre dans le gain.
Un revenu déclaré plus tard, autrement ou dans une société est traité comme un effet sur la recette, net de ce qui est
récupéré ailleurs (la moitié de la réaction des hauts revenus passe à l'IS et au prélèvement forfaitaire, à 48,6~\%,
convention), donné à part et sans rien financer, ce qui minore le gain si l'optimisation consomme des ressources
réelles. Une réaction de calendrier, comme celle des dividendes en 2013, 2018 et 2024, ne compte presque pas.

\paragraph{Les marges.} Les heures réagissent au taux marginal ; la décision de travailler, sous deux Smic, au taux de
prélèvement sur tout le premier salaire. Ni les pensions, ni les allocations chômage, ni l'épargne ne réagissent, faute
de mesure ; l'âge de départ et le cumul emploi-retraite ne sont pas modélisés.

\subsection{Quatre jeux d'hypothèses sur les comportements}\label{ann:meth_jeux}

Les élasticités ont été mesurées face à de petites réformes, d'une année sur l'autre ; nul ne sait comment les
comportements réagiraient à un système entièrement différent. D'où trois jeux, bas, central (prudent, calé sur les
études françaises autant que possible) et haut, plus pour la fusion un jeu de long terme (tableaux~\ref{tab:ircsg_jeux}
et~\ref{tab:meth_jeux}).

\begin{table}[!htbp]\centering\scriptsize\setlength{\tabcolsep}{4pt}
\caption{Les paramètres de comportement de l'immobilier, de l'IS et des revenus du capital}\label{tab:meth_jeux}
\begin{tabularx}{\linewidth}{@{}>{\raggedright\arraybackslash}X>{\raggedright\arraybackslash}p{1.5cm}>{\raggedright\arraybackslash}p{1.5cm}>{\raggedright\arraybackslash}p{1.9cm}>{\raggedright\arraybackslash}p{3.4cm}@{}}
\toprule
Paramètre & Bas & Central & Haut & Appui \\
\midrule
Immobilier : ventes en moins par point de droits & 0 & 6~\% & 15~\% & convention, d'après Bérard et Trannoy (2018), Eerola et al. (2021) \\
Immobilier : capital logement, élasticité au coût d'usage & 0,5 & 1 & 1 (coût d'usage de 3,5~\%) & convention \\
IS : stock d'équipements, élasticité au coût d'usage & 0,5 & 1 & 1,6 & partie IS (H18-H19) ; borne basse sans source \\
Revenus du capital : revenus déclarés, élasticité au taux net & 0,65 & 0,77 & 1 & Lefebvre, Lehmann et Sicsic (2025) ; CAE (2025) \\
Revenus du capital : part de la réaction qui échappe définitivement à l'impôt & 0 & 0 & 1 & non mesurée \\
Épargne : élasticité au rendement net & 0 & 0 & 0 (0,1 à 0,5 en variante) & aucune estimation française \\
\bottomrule
\end{tabularx}
\par\vspace{2pt}{\raggedright\scriptsize Des choix appuyés sur les études, non des estimations. Les jeux bas et haut
combinent toutes les hypothèses basses ou hautes : ce sont des bornes, pas des scénarios.\par}
\end{table}

\paragraph{Le prix d'un milliard.} L'écart de recette entre deux variantes est supposé levé par un point de taux de
l'impôt fusionné pour tous, au prix de 0,04~Md€ de perte sèche par milliard dans le jeu bas, 0,08 dans le central et
0,17 dans le haut. C'est peu parce que les élasticités sont basses : avec 0,3 et un coin de 40 à 50~\%, la formule
usuelle donnerait 0,25 à 0,43 (avis Cochrane).

\subsection{Conventions communes}\label{ann:meth_conv}

\paragraph{Référence.} Le droit en vigueur en novembre 2026, sans le projet de loi de finances pour 2027, donc avec la
fin de la contribution exceptionnelle sur l'IS et l'extinction de la CRDS, que la synthèse affiche à part. Chaque
réforme est calibrée à recette constante avant toute réaction.

\paragraph{Base 2025.} Les chiffrages sont calés sur les prélèvements obligatoires de 2025, dans le sens que précise
l'annexe de chaque partie (pour l'IS, bénéfices de 2025 portés en 2027 ; pour la CEHR, revenus de 2025). Restent en
2024 la répartition entre foyers, le détail des déclarations, l'impôt sur le revenu imputé aux revenus du capital et
les comparaisons de l'OCDE. La synthèse additionne ces montants sans les indexer : l'erreur, de quelques pour cent, est
petite devant les fourchettes.

\paragraph{Actualisation.} 3,25~\% pour l'État (dette de transition, déduction immédiate), 5~\% pour les entreprises,
3~\% réels pour les ménages (taxe foncière capitalisée, report de paiement, droits annualisés), 3,5~\% nominal (2,5 à
5~\%) pour les héritiers, 1 à 3~\% réels pour la valeur actualisée de la hausse des retraités.

\paragraph{Incidence.} L'impôt sur le revenu et la CSG pèsent sur ceux qui les paient ; la hausse de taxe foncière
passe dans les prix et la baisse des droits profite à l'acheteur (cas central) ; l'IS est compté au niveau de la
société.

\subsection{Les calculs de la synthèse}\label{ann:meth_synthese}

\paragraph{Trajectoire.} Soldes des parties : immobilier sur cinq ans (droits divisés par deux dès l'annonce, taux
national montant d'un cinquième par an), IS au cas central (comptes de l'exercice, en caisse un an plus tard),
effacement au jeu central ; un an de retard sur 2027 décalerait la dette sans la changer. Le cas défavorable combine
90~\% d'investissement fait par des sociétés imposables et le jeu bas de l'effacement, le favorable 60~\%, le jeu haut
et la reprise des ventes.

\paragraph{Gain total.} Celui de l'IS est la moyenne des deux régimes d'amortissement actuels à une élasticité de 1,
borné par le dégressif à 0,5 et le linéaire à 1,6. Les bornes de l'immobilier sont les jeux bas et haut du
tableau~\ref{tab:meth_jeux}. Le gain permanent retire de celui de la fusion la part des retraités actuels, selon que
les actifs perçoivent à moitié ou entièrement l'impôt futur sur leur pension (\ref{ann:ircsg}).

\paragraph{Cas-types.} Impôt fusionné des cas-types OpenFisca de la partie IR-CSG ; logement médian de la commune
(Limoges, 171\,800~€ ; Paris, 416\,200~€ ; Lyon, 279\,200~€), sans effet de la taxe sur les prix ni sur les loyers ;
droits annualisés sur dix ans à 3~\% ; PME industrielle de la partie IS ; impôt transmis à 31,4~\% de la plus-value
latente supposée (80\,000~€ pour le couple de retraités, 120\,000~€ pour le cadre, 5~M€ pour le dirigeant, 40~M€ pour
le rentier).

\begin{table}[!htbp]\centering\scriptsize
\caption{Hypothèses propres à la synthèse}\label{tab:meth_hyp}
\begin{tabularx}{\linewidth}{@{}>{\raggedright\arraybackslash}p{5.2cm}>{\raggedright\arraybackslash}p{2.2cm}>{\raggedright\arraybackslash}X@{}}
\toprule
Hypothèse & Où elle joue & Si elle est fausse \\
\midrule
Le 1~\% paie 90~\% de la fin de l'effacement & bilan du haut & À 75~\%, 2,1~Md€ au lieu de 2,5 \\
Le 1~\% paie 0 à 0,5~Md€ du taux national sur le sol des entreprises et les terrains nus & bilan du haut & Non vérifiable :
détention non publiée par niveau de revenu \\
50 à 80~\% de la CEHR portent sur les revenus du capital & bilan du haut & Ne fait que répartir un coût de 1,4~Md€ \\
IR imputé au capital (revenus de 2024) : 10,2~Md€ au lieu de 11,5 en variante & bilan du haut & Majorant ; l'impôt de
l'assurance-vie, non chiffré, réduirait un peu la correction \\
Report de taxe foncière réservé aux 65 ans et plus : moitié des hausses & dette brute & Ordre de grandeur non vérifié ;
majorant de 20~Md€ fin 2036 \\
Taxe sur le sol sans effet sur les loyers & cas-types & Effet théorique, non mesuré \\
\bottomrule
\end{tabularx}
\end{table}

\subsection{Le bilan du 1~\% le plus aisé}\label{ann:meth_haut}

Le tableau~\ref{tab:ens_haut} fait le compte pour le centile des foyers au revenu fiscal de référence le plus élevé.
La fusion supprime la CEHR sur les revenus du barème, ce que la partie IR-CSG ne compte pas. Restent non répartis le
taux national sur le sol des entreprises et des terrains nus (1,4~Md€ net) et la fin du taux réduit de l'IS, qui touche
des holdings patrimoniales et des sociétés de professions libérales.

\begin{table}[!htbp]\centering\scriptsize
\caption{La charge du 1~\% le plus aisé (Md€ par an ; + : paie davantage)}\label{tab:ens_haut}
\begin{tabularx}{\linewidth}{@{}>{\raggedright\arraybackslash}X>{\raggedright\arraybackslash}p{3.2cm}>{\raggedright\arraybackslash}p{2.6cm}>{\raggedright\arraybackslash}p{2.6cm}@{}}
\toprule
 & Convention centrale (fourchette) & IR imputé au capital corrigé & Revenus du capital logés tout en haut \\
\midrule
Fusion IR-CSG & 0,0 & $-$1,1 & $-$3,0 \\
CEHR sur les revenus du barème & $-$0,5 ($-$0,3 à $-$0,7) & $-$0,5 & $-$0,5 \\
Fin de l'effacement (moyenne sur dix ans) & +2,5 (0,7 à 6,8) & +2,5 & +2,5 \\
Immobilier (sol des entreprises et terrains nus) & 0 à +0,5, non compté & id. & id. \\
Fin du taux réduit de l'IS & non chiffré & id. & id. \\
\textbf{Bilan} & \textbf{+2,0 ($-$0,1 à 7,0)} & \textbf{+0,9} & \textbf{$-$1,0} \\
Suppression de la majoration ex-CEHR sur le capital & $-$0,9 & $-$0,9 & $-$0,9 \\
\textbf{Après suppression} & \textbf{+1,1 ($-$1,0 à 6,1)} & \textbf{0,0} & \textbf{$-$1,9} \\
Barème adouci au-delà de 150\,000~€ & $-$1,1, soit 0,0 & --- & --- \\
\bottomrule
\end{tabularx}
\par\vspace{2pt}{\raggedright\scriptsize Hypothèses : tableau~\ref{tab:meth_hyp}. Source :
\texttt{sorties/synthese\_ensemble.txt}, \texttt{controle\_hauts\_purs.txt}.\par}
\end{table}

Le bilan dépend surtout de la répartition par tranche de revenu des revenus au prélèvement forfaitaire, que la DGFiP ne
publie pas. Le modèle IR-CSG les répartit au prorata des autres revenus, les loge d'abord tout en haut (« haut », que
suggère un calage sur l'impôt observé), ou prend la moyenne des deux (« mixte », centrale). Second écart, de même sens :
sur les revenus de 2024, le modèle retire de l'impôt sur le revenu 11,5~Md€ imputés aux revenus du capital, la partie
PFU au plus 10,2, sur l'assiette du prélèvement forfaitaire (79,6~Md€), qui exclut l'assurance-vie et les plus-values
immobilières. Corrigé, l'impôt des revenus du barème est sous-estimé d'environ 1,4~Md€ en 2025 : la fusion rendrait
cette somme, dont 1,1~Md€ au 1~\%, et la recette à remplacer passerait à environ 231,3~Md€ ; tenir inchangée la charge du 1~\% à recette constante
obligerait alors à atteindre 50~\% dès 180\,000~€ et ferait bouger de 540~€ la charge moyenne d'un dixième. Le modèle porte en outre cet
impôt au niveau de 2025 (+7,9~\%), quand la partie PFU le mesure en baisse de 18~\% sur les revenus de 2025 : ces
conventions ne sont pas réconciliées.

La suppression de la majoration ex-CEHR coûte 0,9~Md€ (0,7 à 1,1), un peu moins tant que la contribution
différentielle sur les hauts revenus reprend la baisse des rentiers à son plancher : le rentier des cas-types n'y gagne
rien avant qu'elle disparaisse, 107\,500~€ par an ensuite. Calculée sur la CEHR des revenus de 2024, encaissée en
2025, elle coûterait 1,9~Md€, le 1~\% paierait 0,5~Md€ de plus et la dette de transition culminerait à 35~Md€.

Si le calage confirme la convention centrale, rendre encore environ 1,1~Md€ au 1~\% passerait au mieux par le haut du
barème : le garder jusqu'à 150\,000~€ de revenu imposé (41,5~\%), puis n'atteindre 50~\% qu'à 1,25~M€ au lieu de
231\,000~€. Seul le 1~\% est touché, aucun dixième ne bouge de plus de 360~€, le gain de la fusion monte de 4,7 à
4,9~Md€ et la dette de transition à 39~Md€. Le faire à recette constante ferait payer les autres foyers, dont les
retraités ($+$0,4~Md€) : nous ne le recommandons pas. Si le calage loge les revenus du capital tout en haut, il faut
garder la majoration.

\paragraph{Calcul.} Le modèle IR-CSG est relancé sans modification dans chaque convention, et une fois avec 10,2~Md€
d'impôt imputé, réduit dans la même proportion dans chaque tranche. Le barème adouci ajoute au barème proposé un nœud à
150\,000~€, d'où le taux monte en ligne droite jusqu'à 50~\% à un seuil cherché par dichotomie.

\subsection{Programmes}\label{ann:meth_programmes}

{\small\raggedright Depuis \texttt{fiscalite/} : \texttt{synthese/synthese\_ensemble.py} (sorties
\texttt{synthese/sorties/synthese\_ensemble.txt} et \texttt{.json}), puis \texttt{synthese/figures\_ensemble.py}, une
seconde fois avec \texttt{FIG\_NU=1} ; perte sèche de l'IS : \texttt{synthese/perte\_seche\_autres.py} ; méthode :
\texttt{revue/methode\_v2.md}. Les programmes des parties sont dans leurs annexes.\par}
